\section{Create a map}
We want to create a map of the magnetic field lines using PyOculus. The magnetic configuration considered is described using slab geometry through
\begin{equation}
\psi=\psi_0(x)+\psi_1(x,y,z)\;,
\end{equation}
where $\psi_0$ describes the equilibrium configuration and $\psi_1$ is a perturbation of amplitude $\epsilon$ with poloidal and toroidal mode numbers $m$ and $n$:
\begin{align}
&\psi_0(x)=\frac{\iota'}{2}(x-x_0)^2\,,\\
&\psi_1(x,y,z)=\epsilon(x)\cos\alpha\equiv \epsilon(x)\cos(my-nz)\,.
\end{align}

For a system with $x\in[0,1]$ containing an island at $x_s=0.5$, as a simple model for a perturbation localized around the rational surface, we choose
\begin{equation}\label{pert_ampl}
\epsilon(x)=\delta\cdot x(1-x)\,.
\end{equation}

In PyOculus, all the calculations are performed in cylindrical geometry. For now, we'll consider a simple mapping between the slab and cylindrical geometries, where $x$ is mapped to the radial coordinate $r$, and $y$ and $z$ are mapped to the poloidal angle $\theta$ and toroidal angle $\phi$, respectively.

As for the magnetic configuration itself, there is no specific function in PyOculus that reproduces exactly the case we wanted to study.


